% Source: Gerzensee "Calculus" slides 3-42 (the two firm problems that frame the
% chapter; Definition 3.1 of the derivative; the first-order condition and its
% intuition; the FOC for concave/convex functions and the tangent-line picture;
% Propositions 3.3-3.6 (linearity, product, quotient with two derivations, chain
% rule); the derivative of an inverse function; price-taking and price-setting
% firms, elasticity and the markup, Exercise 11.14; the Taylor coefficient
% derivation; the exp and log(1+x) expansions; growth accounting, growth rates,
% the Fisher equation, log-linearization; the risk premium and the variance
% identity).
% Supplementary: Fukuda (2026), Ch. 3 (single-variable calculus) and Sec. 11.2,
% 11.3, 11.7.
% No external reference is cited in this chapter: every statement is drawn
% from the Gerzensee slides and the supplementary lecture notes listed above.

\chapter{Differential Calculus in One Variable}
\label{ch:calculus1}

\section{Motivation: The Two Firm Problems}\label{sec:c1-motivation}

Two problems organise this chapter
\autocite[slide 3]{fukuda2026calculus}. A price-taking firm solves
\begin{equation}\label{eq:price-taker}
	\max_{q\in[0,\infty)}\ pq-c(q),
\end{equation}
and a price-setting firm facing inverse demand $P$ solves
\begin{equation}\label{eq:price-setter}
	\max_{q\in[0,\infty)}\ P(q)q-c(q).
\end{equation}

\begin{assumption}[Standing assumptions on the firm]\label{ass:firm}
	The cost function $c\colon[0,\infty)\to\R$ is differentiable, strictly
	increasing and strictly convex, with $c(0)=c'(0)=0$ and
	$\lim_{q\to\infty}c(q)=\lim_{q\to\infty}c'(q)=\infty$. The revenue function
	$q\mapsto P(q)q$ is concave and differentiable; the leading case is
	$P(q)=a-bq$ with $a,b>0$.
\end{assumption}

Each assumption is used at an identifiable point. Differentiability makes the
first-order condition available; strict convexity of $c$ and concavity of
revenue make the first-order condition sufficient
(\autoref{thm:foc-concave}) and the solution unique; $c'(0)=0$ with
$c'(q)\to\infty$ places the optimum in the interior of $[0,\infty)$, where
\autoref{thm:fermat} applies; and $c(0)=0$ makes the value of the problem the
profit rather than profit net of a fixed cost.

\section{The Derivative}\label{sec:c1-derivative}

\begin{definition}[Derivative]\label{def:derivative}
	Let $f\colon(a,b)\to\R$ with $-\infty\leq a<b\leq\infty$ and let
	$x_0\in(a,b)$. The function $f$ is \dfnas{differentiable at}{differentiability}
	$x_0$ if
	\[
		f'(x_0)\eqdef\lim_{h\to0}\frac{f(x_0+h)-f(x_0)}{h}
	\]
	exists and is finite; the limit is the \dfn{derivative} of $f$ at $x_0$, also
	written $Df(x_0)$ or $\deriv{f}{x}(x_0)$. It is the slope of the tangent line to
	the graph of $f$ at $\bigl(x_0,f(x_0)\bigr)$.
	\nidx{f'(x)}{$f'(x)$, derivative}
\end{definition}

This is Definition~3.1 of \textcite{fukuda2026notes}
\autocite[slide 4]{fukuda2026calculus}. That $x_0$ is required to be
\emph{interior} to the domain is part of the definition: the two-sided limit
must exist, so the difference quotient must be defined for $h$ of both signs.

\begin{proposition}[Equivalent formulation and continuity]\label{prop:diff-continuous}
	Let $f\colon(a,b)\to\R$ and $x_0\in(a,b)$. Then $f$ is differentiable at $x_0$
	with derivative $L$ if and only if
	\begin{equation}\label{eq:linear-approx}
		f(x_0+h)=f(x_0)+Lh+r(h),\qquad\text{where}\quad\lim_{h\to0}\frac{r(h)}{h}=0 .
	\end{equation}
	In particular, differentiability at $x_0$ implies continuity at $x_0$.
\end{proposition}

\begin{proof}
	Define $r(h)\eqdef f(x_0+h)-f(x_0)-Lh$. Then $r(h)/h$ equals the difference
	quotient minus $L$, so $r(h)/h\to0$ if and only if the difference quotient tends
	to $L$. For continuity, \eqref{eq:linear-approx} gives
	$f(x_0+h)-f(x_0)=Lh+r(h)\to0$ as $h\to0$, since $r(h)=h\cdot(r(h)/h)\to0\cdot0$.
\end{proof}

Formulation \eqref{eq:linear-approx} is the one that generalises: it says that
$f$ is well approximated near $x_0$ by an affine function, with an error small
relative to $h$. In one variable the approximating linear map is multiplication
by a number; in $\R^n$ it is a linear map, and \autoref{def:total-derivative}
takes \eqref{eq:linear-approx} verbatim as the definition. The converse of the
last sentence fails: $f(x)=\abs{x}$ is continuous at $0$ and not differentiable
there, since the difference quotient equals $\sgn(h)$.

\subsection{The first-order condition}

\begin{theorem}[Fermat]\label{thm:fermat}
	Let $A\subseteq\R$, let $x_0$ be an interior point of $A$ at which
	$f\colon A\to\R$ is differentiable, and suppose $x_0$ is a local maximiser or a
	local minimiser of $f$ on $A$. Then $f'(x_0)=0$.
\end{theorem}

\begin{proof}
	Suppose $x_0$ is a local maximiser, so there is $\delta>0$ with
	$\ball{\delta}{x_0}\subseteq A$ and $f(x)\leq f(x_0)$ for
	$x\in\ball{\delta}{x_0}$. For $0<h<\delta$,
	\[
		\frac{f(x_0+h)-f(x_0)}{h}\leq0,
	\]
	because the numerator is non-positive and the denominator positive; letting
	$h\downarrow0$ gives $f'(x_0)\leq0$. For $-\delta<h<0$ the denominator is
	negative and the quotient is non-negative; letting $h\uparrow0$ gives
	$f'(x_0)\geq0$. Since $f$ is differentiable at $x_0$ the two one-sided limits
	coincide with $f'(x_0)$, so $f'(x_0)=0$. The minimiser case follows by applying
	this to $-f$.
\end{proof}

A point with $f'(x_0)=0$ is a \dfnas{stationary point}{stationary point}, or
critical point, of $f$ \autocite[slide 5]{fukuda2026calculus}.

\begin{remark}[Three ways the conclusion fails]\label{rem:foc-failures}
	Each hypothesis of \autoref{thm:fermat} is indispensable, and each failure has
	an economic counterpart.
	\begin{enumerate}[(i)]
		\item \emph{Not interior.} $f(x)=x$ on $[0,1]$ has its maximum at $x=1$ with
		      $f'(1)=1\neq0$. The proof used $h$ of both signs, and at a boundary point
		      only one sign is admissible; the conclusion weakens to the inequality
		      $f'(1)\geq0$. In economics this is the corner solution, and the
		      corresponding first-order conditions are the Kuhn--Tucker inequalities of
		      \autoref{ch:optimization}.
		\item \emph{Not differentiable.} $f(x)=-\abs{x}$ has a global maximum at
		      $x=0$ where no derivative exists. Non-differentiable kinks arise in
		      economics from piecewise-linear tax schedules and from Leontief technology.
		\item \emph{The condition is necessary, not sufficient.} $f(x)=x^3$ has
		      $f'(0)=0$ and no local extremum at $0$. Sufficiency requires either a
		      second-order condition or a global structural hypothesis such as concavity.
	\end{enumerate}
\end{remark}

\subsection{Concavity and sufficiency}

\begin{definition}[Concave and convex functions of one variable]\label{def:concave-1d}
	A function $f\colon I\to\R$ on an interval $I\subseteq\R$ is
	\dfnas{concave}{concave function} if
	\[
		f\bigl(\lambda x+(1-\lambda)y\bigr)\geq\lambda f(x)+(1-\lambda)f(y)
		\qquad\text{for all }x,y\in I,\ \lambda\in[0,1],
	\]
	\dfn{strictly concave} if the inequality is strict whenever $x\neq y$ and
	$\lambda\in(0,1)$, and \dfnas{convex}{convex function} (respectively strictly
	convex) if $-f$ is concave (respectively strictly concave). The general theory,
	in $\R^n$, is \autoref{ch:convex}.
\end{definition}

\begin{theorem}[Characterisation of differentiable concavity]\label{thm:concave-tangent}
	Let $f\colon(a,b)\to\R$ be differentiable. The following are equivalent.
	\begin{enumerate}[(i)]
		\item $f$ is concave on $(a,b)$.
		\item $f(x)\leq f(x_0)+f'(x_0)(x-x_0)$ for all $x,x_0\in(a,b)$.
		\item $f'$ is non-increasing on $(a,b)$.
	\end{enumerate}
	If moreover $f$ is twice differentiable, each is equivalent to $f''\leq0$ on
	$(a,b)$. The strict versions correspond to strict inequality in (ii) for
	$x\neq x_0$ and to strictly decreasing $f'$; $f''<0$ is sufficient but not
	necessary for strict concavity, as $f(x)=-x^4$ shows at $x=0$.
\end{theorem}

\begin{proof}
	(i)$\Rightarrow$(ii). For $\lambda\in(0,1]$, concavity gives
	$f\bigl(x_0+\lambda(x-x_0)\bigr)\geq\lambda f(x)+(1-\lambda)f(x_0)$, hence
	\[
		\frac{f\bigl(x_0+\lambda(x-x_0)\bigr)-f(x_0)}{\lambda}\geq f(x)-f(x_0).
	\]
	Letting $\lambda\downarrow0$, the left side tends to $f'(x_0)(x-x_0)$ by
	\autoref{def:derivative} applied with $h=\lambda(x-x_0)$.

	(ii)$\Rightarrow$(iii). Apply (ii) twice, at $x_0$ evaluated at $x$ and at $x$
	evaluated at $x_0$, and add:
	$0\leq\bigl(f'(x_0)-f'(x)\bigr)(x-x_0)$. For $x>x_0$ this forces
	$f'(x)\leq f'(x_0)$.

	(iii)$\Rightarrow$(i). Let $x<y$ and $z=\lambda x+(1-\lambda)y$ with
	$\lambda\in(0,1)$. By the mean value theorem (\autoref{thm:mvt}) there are
	$\xi_1\in(x,z)$ and $\xi_2\in(z,y)$ with
	$f(z)-f(x)=f'(\xi_1)(z-x)$ and $f(y)-f(z)=f'(\xi_2)(y-z)$. Since
	$\xi_1<\xi_2$, (iii) gives $f'(\xi_1)\geq f'(\xi_2)$, so
	\[
		\frac{f(z)-f(x)}{z-x}\geq\frac{f(y)-f(z)}{y-z} .
	\]
	Substituting $z-x=(1-\lambda)(y-x)$ and $y-z=\lambda(y-x)$ and rearranging gives
	$f(z)\geq\lambda f(x)+(1-\lambda)f(y)$.

	The twice-differentiable statement is (iii) together with the fact that a
	differentiable function is non-increasing exactly when its derivative is
	non-positive, itself a consequence of \autoref{thm:mvt}.
\end{proof}

\begin{theorem}[The first-order condition is sufficient under concavity]
	\label{thm:foc-concave}
	Let $f\colon(a,b)\to\R$ be differentiable and concave. Then $x_0\in(a,b)$ is a
	global maximiser of $f$ on $(a,b)$ if and only if $f'(x_0)=0$. If $f$ is
	convex, the same statement holds with ``maximiser'' replaced by ``minimiser''.
	If $f$ is strictly concave, the maximiser is unique when it exists.
\end{theorem}

\begin{proof}
	Necessity is \autoref{thm:fermat}. For sufficiency, suppose $f'(x_0)=0$. By
	\autoref{thm:concave-tangent}(ii), $f(x)\leq f(x_0)+f'(x_0)(x-x_0)=f(x_0)$ for
	every $x\in(a,b)$. Uniqueness under strict concavity: if $x_0\neq x_1$ were both
	maximisers with common value $M$, then
	$f\bigl(\tfrac12x_0+\tfrac12x_1\bigr)>\tfrac12M+\tfrac12M=M$, contradicting
	maximality.
\end{proof}

This is the content of \autocite[slides 7--8]{fukuda2026calculus}. The geometric
statement --- the tangent line at any point lies above the graph, so a horizontal
tangent lies above the whole graph --- is the proof, once
\autoref{thm:concave-tangent} makes ``lies above'' precise.

\subsection{Mean value theorems}

\begin{theorem}[Rolle]\label{thm:rolle}
	Let $f\colon[a,b]\to\R$ be continuous on $[a,b]$ and differentiable on
	$(a,b)$, with $f(a)=f(b)$. Then $f'(\xi)=0$ for some $\xi\in(a,b)$.
\end{theorem}

\begin{proof}
	$[a,b]$ is compact and $f$ continuous, so by
	\autoref{cor:weierstrass-continuous} $f$ attains a maximum at some
	$x_{\max}$ and a minimum at some $x_{\min}$ in $[a,b]$. If both are endpoints
	then $f$ is constant, since $f(a)=f(b)$ forces $\max f=\min f$, and any
	$\xi\in(a,b)$ works. Otherwise one of them is interior, and
	\autoref{thm:fermat} applies there.
\end{proof}

\begin{theorem}[Mean value theorem]\label{thm:mvt}
	Let $f\colon[a,b]\to\R$ be continuous on $[a,b]$ and differentiable on
	$(a,b)$. Then there is $\xi\in(a,b)$ with
	\[
		f(b)-f(a)=f'(\xi)(b-a).
	\]
\end{theorem}

\begin{proof}
	Apply \autoref{thm:rolle} to
	$g(x)\eqdef f(x)-f(a)-\frac{f(b)-f(a)}{b-a}(x-a)$, which is continuous on
	$[a,b]$, differentiable on $(a,b)$, and satisfies $g(a)=g(b)=0$. Then
	$0=g'(\xi)=f'(\xi)-\frac{f(b)-f(a)}{b-a}$.
\end{proof}

The mean value theorem is the device by which local information about $f'$ is
converted into global information about $f$: it is what makes ``$f'\geq0$ on an
interval'' imply ``$f$ is non-decreasing there'', a step used without comment
whenever a comparative-static sign is read off a derivative.

\section{Rules of Differentiation}\label{sec:c1-rules}

\begin{proposition}[Linearity]\label{prop:deriv-linear}
	If $f,g\colon(a,b)\to\R$ are differentiable at $x$ and $\alpha,\beta\in\R$, then
	$\alpha f+\beta g$ is differentiable at $x$ and
	$(\alpha f+\beta g)'(x)=\alpha f'(x)+\beta g'(x)$.
\end{proposition}

\begin{proof}
	The difference quotient of $\alpha f+\beta g$ equals
	$\alpha$ times that of $f$ plus $\beta$ times that of $g$, and limits are linear
	by \autoref{prop:continuity-algebra}(i).
\end{proof}

\begin{proposition}[Product rule]\label{prop:deriv-product}
	If $f,g$ are differentiable at $x$, so is $fg$, and
	$(fg)'(x)=f'(x)g(x)+f(x)g'(x)$.
\end{proposition}

\begin{proof}
	Add and subtract $f(x+h)g(x)$:
	\[
		\frac{f(x+h)g(x+h)-f(x)g(x)}{h}
		=\frac{f(x+h)-f(x)}{h}\,g(x)+f(x+h)\,\frac{g(x+h)-g(x)}{h}.
	\]
	As $h\to0$ the first quotient tends to $f'(x)$ and the second to $g'(x)$, while
	$f(x+h)\to f(x)$ by \autoref{prop:diff-continuous}.
\end{proof}

The choice of which factor to hold fixed is the substance of the proof: the
decomposition corresponds to splitting the area increment of a rectangle of sides
$f$ and $g$ into two strips, and the omitted corner term
$(f(x+h)-f(x))(g(x+h)-g(x))$ is of order $h^2$
\autocite[slide 11]{fukuda2026calculus}.

\begin{proposition}[Quotient rule]\label{prop:deriv-quotient}
	If $f,g$ are differentiable at $x$ and $g(x)\neq0$, then $f/g$ is
	differentiable at $x$ and
	\[
		\left(\frac{f}{g}\right)'(x)=\frac{f'(x)g(x)-f(x)g'(x)}{g(x)^2}.
	\]
\end{proposition}

\begin{proof}
	Since $g$ is continuous at $x$ and $g(x)\neq0$, $g$ is non-zero on a
	neighbourhood of $x$, so $f/g$ is defined there. Write $h\eqdef f/g$, so
	$f=gh$. By \autoref{prop:deriv-product}, $f'=g'h+gh'$, whence
	\[
		h'=\frac{f'-g'h}{g}=\frac{f'-g'f/g}{g}=\frac{f'g-fg'}{g^2}.
	\]
	The step presumes $h$ differentiable; to avoid that, compute directly
	\[
		\frac{1}{h}\left(\frac{1}{g(x+h)}-\frac{1}{g(x)}\right)
		=-\frac{1}{g(x+h)g(x)}\cdot\frac{g(x+h)-g(x)}{h}
		\;\longrightarrow\;-\frac{g'(x)}{g(x)^2},
	\]
	which establishes $(1/g)'=-g'/g^2$, and then apply
	\autoref{prop:deriv-product} to $f\cdot(1/g)$.
\end{proof}

Both derivations are those of \autocite[slides 13--14, 16]{fukuda2026calculus};
the second is logically self-contained, whereas the first assumes what it
computes and is a mnemonic rather than a proof. The distinction is recorded
because the first is frequently presented as though it were the second.

\begin{proposition}[Chain rule]\label{prop:chain-rule-1d}
	Let $f\colon(a,b)\to\R$ and $g\colon(c,d)\to\R$ with $f\bigl((a,b)\bigr)
		\subseteq(c,d)$. If $f$ is differentiable at $x_0$ and $g$ is differentiable at
	$f(x_0)$, then $g\circ f$ is differentiable at $x_0$ and
	\[
		(g\circ f)'(x_0)=g'\bigl(f(x_0)\bigr)f'(x_0).
	\]
\end{proposition}

\begin{proof}
	Write $y_0=f(x_0)$ and define
	\[
		\varphi(y)\eqdef
		\begin{cases}
			\dfrac{g(y)-g(y_0)}{y-y_0}, & y\neq y_0, \\[1.5ex]
			g'(y_0),                    & y=y_0 .
		\end{cases}
	\]
	Then $\varphi$ is continuous at $y_0$ by differentiability of $g$, and
	$g(y)-g(y_0)=\varphi(y)(y-y_0)$ for \emph{all} $y$, including $y=y_0$. Hence for
	$h\neq0$,
	\[
		\frac{g(f(x_0+h))-g(f(x_0))}{h}
		=\varphi\bigl(f(x_0+h)\bigr)\cdot\frac{f(x_0+h)-f(x_0)}{h}.
	\]
	As $h\to0$, $f(x_0+h)\to y_0$ by continuity of $f$, so the first factor tends to
	$\varphi(y_0)=g'(y_0)$ by \autoref{prop:sequential-continuity}, and the second
	to $f'(x_0)$.
\end{proof}

The auxiliary function $\varphi$ exists to handle the case $f(x_0+h)=f(x_0)$, at
which the naive multiply-and-divide argument would divide by zero. The naive
argument is the one behind the suggestive notation
$\deriv{z}{x}=\deriv{z}{y}\deriv{y}{x}$ of
\autocite[slide 18]{fukuda2026calculus}; it gives the right answer and is not a
proof.

\begin{proposition}[Derivative of an inverse function]\label{prop:inverse-derivative}
	Let $f\colon(a,b)\to\R$ be continuous and strictly monotone, hence injective
	with continuous inverse $f^{-1}$ on the interval $f\bigl((a,b)\bigr)$. If $f$
	is differentiable at $x_0$ with $f'(x_0)\neq0$, then $f^{-1}$ is differentiable
	at $y_0=f(x_0)$ and
	\[
		\bigl(f^{-1}\bigr)'(y_0)=\frac{1}{f'\bigl(f^{-1}(y_0)\bigr)}
		=\frac{1}{f'(x_0)} .
	\]
\end{proposition}

\begin{proof}
	Let $y_k\to y_0$ with $y_k\neq y_0$ and set $x_k\eqdef f^{-1}(y_k)$, so
	$x_k\neq x_0$ and $x_k\to x_0$ by continuity of $f^{-1}$. Then
	\[
		\frac{f^{-1}(y_k)-f^{-1}(y_0)}{y_k-y_0}
		=\frac{x_k-x_0}{f(x_k)-f(x_0)}
		=\left(\frac{f(x_k)-f(x_0)}{x_k-x_0}\right)^{-1}
		\longrightarrow\frac{1}{f'(x_0)},
	\]
	the last step by $f'(x_0)\neq0$ and continuity of $t\mapsto1/t$ away from $0$.
	\autoref{prop:sequential-continuity} converts the sequential statement into the
	limit.
\end{proof}

The hypothesis $f'(x_0)\neq0$ cannot be dropped: $f(x)=x^3$ is a strictly
increasing bijection of $\R$ with $f'(0)=0$, and $f^{-1}(y)=y^{1/3}$ has a
vertical tangent at $0$. This is the one-dimensional case of the inverse function
theorem, and the condition $f'(x_0)\neq0$ becomes invertibility of the Jacobian
in \autoref{thm:inverse-function}. Economically it is the statement that a demand
function $D$ and an inverse demand function $P$ have reciprocal slopes,
$P'(q)=1/D'(p)$ at $q=D(p)$ \autocite[slides 19--20]{fukuda2026calculus}.

\section{Taylor Expansion}\label{sec:c1-taylor}

The problem is to approximate $f$ near $x_0$ by a polynomial in $h$,
\begin{equation}\label{eq:taylor-ansatz}
	f(x_0+h)\approx\alpha_0+\alpha_1h+\alpha_2h^2+\cdots+\alpha_nh^n .
\end{equation}
Setting $h=0$ gives $\alpha_0=f(x_0)$. Differentiating \eqref{eq:taylor-ansatz}
$k$ times with respect to $h$ and again setting $h=0$ kills every term of degree
below $k$ and every term of degree above $k$, leaving
$f^{(k)}(x_0)=k!\,\alpha_k$, so
\begin{equation}\label{eq:taylor-coefficients}
	\alpha_k=\frac{f^{(k)}(x_0)}{k!},\qquad k=0,1,\dots,n .
\end{equation}
This is the derivation of \autocite[slides 28--31]{fukuda2026calculus}. It
identifies the coefficients on the assumption that an expansion of the stated
form exists; it does not establish that one does, nor how large the error is.
Both are the content of the following theorem.

\begin{theorem}[Taylor's theorem with Lagrange remainder]\label{thm:taylor-1d}
	Let $n\in\N_0$, let $I\subseteq\R$ be an open interval containing $x_0$, and let
	$f\colon I\to\R$ be $n+1$ times differentiable on $I$. Define the
	\dfnas{Taylor polynomial}{Taylor polynomial} of order $n$,
	\[
		T_nf(x_0;h)\eqdef\sum_{k=0}^{n}\frac{f^{(k)}(x_0)}{k!}h^k .
	\]
	Then for every $h$ with $x_0+h\in I$ there exists $\xi$ strictly between $x_0$
	and $x_0+h$ (or $\xi=x_0$ if $h=0$) such that
	\begin{equation}\label{eq:lagrange-remainder}
		f(x_0+h)=T_nf(x_0;h)+\frac{f^{(n+1)}(\xi)}{(n+1)!}\,h^{n+1}.
	\end{equation}
\end{theorem}

\begin{proof}
	Fix $h\neq0$ and let $J$ be the closed interval with endpoints $x_0$ and
	$x_0+h$. Define $M$ by
	$f(x_0+h)=T_nf(x_0;h)+M h^{n+1}$, which determines $M$ uniquely since
	$h\neq0$, and set
	\[
		g(t)\eqdef f(t)-\sum_{k=0}^{n}\frac{f^{(k)}(x_0)}{k!}(t-x_0)^k-M(t-x_0)^{n+1},
		\qquad t\in J .
	\]
	By construction $g^{(k)}(x_0)=0$ for $k=0,1,\dots,n$, and $g(x_0+h)=0$ by the
	choice of $M$. Since $g(x_0)=g(x_0+h)=0$, \autoref{thm:rolle} gives
	$\xi_1$ strictly between them with $g'(\xi_1)=0$. Since also $g'(x_0)=0$,
	Rolle applied to $g'$ on the interval with endpoints $x_0$ and $\xi_1$ gives
	$\xi_2$ with $g''(\xi_2)=0$. Iterating $n+1$ times produces
	$\xi\eqdef\xi_{n+1}$ strictly between $x_0$ and $x_0+h$ with
	$g^{(n+1)}(\xi)=0$. But $g^{(n+1)}(t)=f^{(n+1)}(t)-(n+1)!\,M$, so
	$M=f^{(n+1)}(\xi)/(n+1)!$, which is \eqref{eq:lagrange-remainder}.
\end{proof}

\begin{corollary}[Peano remainder]\label{cor:peano}
	If $f$ is $n$ times differentiable at $x_0$, then
	$f(x_0+h)=T_nf(x_0;h)+o(h^n)$ as $h\to0$; that is, the error divided by $h^n$
	tends to $0$.
\end{corollary}

The two remainder forms answer different questions. The Lagrange form
\eqref{eq:lagrange-remainder} is an exact identity with an unknown evaluation
point, and it yields \emph{numerical} error bounds once $f^{(n+1)}$ is bounded on
the relevant interval. The Peano form is an asymptotic statement with no
evaluation point, and it is the one used when the argument is a limit as a
parameter goes to zero --- which is exactly the situation in
\autoref{sec:c1-risk} below.

\begin{eg}[Exponential and logarithm]\label{eg:taylor-exp-log}
	For $f(x)=e^x$ every derivative is $e^x$ and $f^{(k)}(0)=1$, so
	\[
		e^x=1+x+\tfrac12x^2+\tfrac16x^3+\cdots+\tfrac{1}{n!}x^n
		+\frac{e^{\xi}}{(n+1)!}x^{n+1}.
	\]
	For $f(x)=\log(1+x)$ on $(-1,\infty)$, $f'(x)=(1+x)^{-1}$,
	$f''(x)=-(1+x)^{-2}$, $f'''(x)=2(1+x)^{-3}$, so $f^{(k)}(0)=(-1)^{k-1}(k-1)!$
	and
	\[
		\log(1+x)=x-\tfrac12x^2+\tfrac13x^3-\dots
		+\frac{(-1)^{n-1}}{n}x^{n}+\frac{(-1)^{n}}{n+1}\frac{x^{n+1}}{(1+\xi)^{n+1}}.
	\]
	The first-order truncations $e^x\approx1+x$ and $\log(1+x)\approx x$ carry an
	error of order $x^2$, with sign determined by the second derivative: the
	exponential approximation understates and the logarithmic one overstates
	\autocite[slides 32--36]{fukuda2026calculus}.
\end{eg}

\subsection{Four uses of the logarithmic approximation}

\begin{eg}[Growth accounting]\label{eg:growth-accounting}
	Since $(\log f)'=f'/f$ by \autoref{prop:chain-rule-1d}, taking logarithms of
	the Cobb--Douglas aggregate $Y(t)=A(t)K(t)^\alpha L(t)^{1-\alpha}$ and
	differentiating in $t$ gives the exact identity
	\[
		\frac{Y'(t)}{Y(t)}=\frac{A'(t)}{A(t)}+\alpha\frac{K'(t)}{K(t)}
		+(1-\alpha)\frac{L'(t)}{L(t)} .
	\]
	No approximation is involved: the growth rate of a Cobb--Douglas aggregate is
	exactly the factor-share-weighted sum of the growth rates of its arguments, and
	the residual $A'/A$ is what growth accounting measures
	\autocite[slide 37]{fukuda2026calculus}.
\end{eg}

\begin{eg}[Log differences as growth rates]\label{eg:log-growth}
	For a positive sequence $(Y_t)$ with growth rate
	$r_t=(Y_t-Y_{t-1})/Y_{t-1}$,
	\[
		\Delta\log Y_t=\log Y_t-\log Y_{t-1}=\log\!\left(\frac{Y_t}{Y_{t-1}}\right)
		=\log(1+r_t)\approx r_t .
	\]
	By \autoref{eg:taylor-exp-log} the error is $-r_t^2/2+O(r_t^3)$, so log
	differences \emph{understate} growth, by about half a percentage point at
	$r_t=10\%$ and by five at $r_t=30\%$. The approximation is a good one at
	quarterly frequencies and a poor one for hyperinflations
	\autocite[slide 38]{fukuda2026calculus}.
\end{eg}

\begin{eg}[Fisher equation]\label{eg:fisher}
	The exact relation between the nominal rate $i$, the inflation rate $\pi$ and
	the real rate $r$ is $(1+i)=(1+\pi)(1+r)$. Taking logarithms turns the product
	into a sum exactly, $\log(1+i)=\log(1+\pi)+\log(1+r)$, and applying
	$\log(1+x)\approx x$ three times gives the Fisher equation
	$i\approx\pi+r$. The neglected term is $\pi r$, which is second order in the two
	rates \autocite[slide 39]{fukuda2026calculus}.
\end{eg}

\begin{eg}[Log-linearisation]\label{eg:log-linearization}
	First-order Taylor expansion gives
	$f(x)\approx f(x_0)+f'(x_0)(x-x_0)$. Writing the absolute change as $x_0$ times
	the proportional change and using $\log x-\log x_0\approx(x-x_0)/x_0$,
	\[
		f(x)\approx f(x_0)+f'(x_0)x_0\cdot\bigl(\log x-\log x_0\bigr).
	\]
	The coefficient $f'(x_0)x_0$ is the semi-elasticity of $f$ at $x_0$. This is the
	step by which a nonlinear dynamic model is converted into a linear system in log
	deviations from steady state, whose stability is then read off the eigenvalues
	as in \autoref{prop:linear-stability} \autocite[slide 40]{fukuda2026calculus}.
\end{eg}

\section{Economic Application: The Firm's Problem}\label{sec:c1-firm}

\subsection{The price-taking firm}

\begin{proposition}\label{prop:price-taker}
	Under \autoref{ass:firm} and for every $p>0$, problem \eqref{eq:price-taker}
	has a unique solution $q^\ast(p)>0$, characterised by
	\begin{equation}\label{eq:mc-equals-price}
		c'(q^\ast)=p,\qquad\text{equivalently}\qquad q^\ast(p)=(c')^{-1}(p).
	\end{equation}
	The supply function $q^\ast$ is strictly increasing and differentiable with
	$\bigl(q^\ast\bigr)'(p)=1/c''\bigl(q^\ast(p)\bigr)>0$ whenever $c$ is twice
	differentiable with $c''>0$.
\end{proposition}

\begin{proof}
	The objective $\pi(q)=pq-c(q)$ is strictly concave, being linear minus strictly
	convex, and differentiable, with $\pi'(q)=p-c'(q)$. Since $c'(0)=0<p$ and
	$c'(q)\to\infty$, the continuous strictly increasing function $c'$ takes the
	value $p$ at exactly one point $q^\ast>0$, by the intermediate value theorem
	(\autoref{cor:ivt}) and injectivity. At that point $\pi'=0$, and
	\autoref{thm:foc-concave} makes it the unique global maximiser on $(0,\infty)$;
	comparison with $\pi(0)=0<\pi(q^\ast)$ rules out the boundary. The supply
	function is the inverse of $c'$, and its derivative is
	\autoref{prop:inverse-derivative}.
\end{proof}

The economic content of \eqref{eq:mc-equals-price} is that the firm equates
marginal cost to price. The comparative static
$(q^\ast)'=1/c''>0$ makes precise the sense in which the slope of the supply
curve is the reciprocal of the curvature of the cost function: a firm with nearly
linear costs responds strongly to price, one with sharply increasing marginal
cost hardly at all.

\subsection{The price-setting firm and the markup}

\begin{proposition}\label{prop:markup}
	Under \autoref{ass:firm}, an interior solution $q^\ast>0$ of
	\eqref{eq:price-setter} satisfies
	\begin{equation}\label{eq:lerner}
		P(q^\ast)\left(1+\frac{1}{\eta}\right)=c'(q^\ast),
		\qquad\text{equivalently}\qquad
		P(q^\ast)=\underbrace{\frac{1}{1+\eta^{-1}}}_{\text{markup}}\;c'(q^\ast),
	\end{equation}
	where
	\[
		\eta\eqdef\restr{\deriv{q}{p}\frac{p}{q}}{q=q^\ast}<0
	\]
	is the price elasticity of demand. If $\eta<-1$ the markup exceeds $1$ and
	$P(q^\ast)>c'(q^\ast)$.
\end{proposition}

\begin{proof}
	The objective is $\pi(q)=P(q)q-c(q)$, concave by \autoref{ass:firm}. Its
	derivative, by \autoref{prop:deriv-product}, is
	$\pi'(q)=P'(q)q+P(q)-c'(q)$, so \autoref{thm:fermat} gives
	\begin{equation}\label{eq:mr-equals-mc}
		\underbrace{P'(q^\ast)q^\ast+P(q^\ast)}_{\text{marginal revenue}}
		=\underbrace{c'(q^\ast)}_{\text{marginal cost}} .
	\end{equation}
	Factoring $P(q^\ast)>0$ out of the left side,
	\[
		P(q^\ast)\left(1+\frac{P'(q^\ast)q^\ast}{P(q^\ast)}\right)=c'(q^\ast).
	\]
	By \autoref{prop:inverse-derivative} applied to the demand--inverse-demand
	pair, $P'(q^\ast)=1/D'(p^\ast)$ at $p^\ast=P(q^\ast)$, so
	\[
		\frac{P'(q^\ast)q^\ast}{P(q^\ast)}
		=\deriv{p}{q}\frac{q}{p}
		=\left(\deriv{q}{p}\frac{p}{q}\right)^{-1}=\frac{1}{\eta},
	\]
	which is \eqref{eq:lerner}. Solving for $P(q^\ast)$ gives the markup form, and
	$\eta<-1$ makes $1+\eta^{-1}\in(0,1)$, hence the markup exceeds $1$.
\end{proof}

Three observations. First, \eqref{eq:lerner} is the Lerner condition, usually
written $\bigl(P-c'\bigr)/P=-1/\eta$; the two forms are algebraically identical.
Second, the derivation uses the inverse function rule twice over --- once to
convert $P'$ into $1/D'$ and once implicitly in reading the elasticity in either
direction --- so \autoref{prop:inverse-derivative} is not a technicality here but
the substance. Third, profit maximisation requires $\eta<-1$: if demand were
inelastic, marginal revenue $P(1+\eta^{-1})$ would be negative while marginal
cost is positive, and \eqref{eq:mr-equals-mc} would have no solution. A monopolist
never operates on the inelastic portion of its demand curve, and this is a
theorem about \eqref{eq:mr-equals-mc} rather than an assumption.

\begin{eg}[The dual formulation]\label{eg:markup-price-choice}
	Choosing the price rather than the quantity, the firm solves
	$\max_{p\geq0}\ pD(p)-c\bigl(D(p)\bigr)$. The first-order condition is
	$D(p^\ast)+p^\ast D'(p^\ast)-c'\bigl(D(p^\ast)\bigr)D'(p^\ast)=0$; dividing by
	$D'(p^\ast)<0$ and rearranging gives
	\[
		p^\ast=\left(\frac{1}{1+\eta^{-1}}\right)c'\bigl(D(p^\ast)\bigr),
		\qquad
		\eta=\restr{\frac{D'(p)p}{D(p)}}{p=p^\ast},
	\]
	the same condition as \eqref{eq:lerner}. That the two formulations agree is a
	consequence of the chain rule and the strict monotonicity of $D$, which makes
	the change of variable $q=D(p)$ a bijection between the two feasible sets. This
	is Exercise~11.14 of \textcite{fukuda2026notes}
	\autocite[slide 25]{fukuda2026calculus}.
\end{eg}

\section{Economic Application: The Risk Premium}\label{sec:c1-risk}

\paragraph{Environment.} A decision maker with wealth $w$ and utility
$u\colon\R\to\R$, strictly increasing and twice continuously differentiable,
faces a gamble $\tilde\varepsilon$ with $\E[\tilde\varepsilon]=0$ and
$\Var[\tilde\varepsilon]=\sigma^2$. The \dfnas{risk premium}{risk premium}
$\pi$ is the sure amount whose surrender leaves the decision maker indifferent to
bearing the gamble:
\begin{equation}\label{eq:risk-premium-def}
	u(w-\pi)=\E\bigl[u(w+\tilde\varepsilon)\bigr].
\end{equation}

The slides derive the approximation by expanding each side and matching
\autocite[slides 26--27, 41]{fukuda2026calculus}. That derivation is correct in
its conclusion and incomplete as an argument: it expands the left side to first
order and the right to second order without saying why the orders may differ, and
it does not say in what limit the approximation is valid. Both gaps are closed by
scaling the gamble.

\begin{theorem}[Arrow--Pratt]\label{thm:arrow-pratt}
	Let $u$ be three times continuously differentiable with $u'>0$ on a
	neighbourhood of $w$, and let $\tilde z$ be a bounded random variable with
	$\E[\tilde z]=0$ and $\Var[\tilde z]=1$. For $t>0$ let $\pi(t)$ solve
	\eqref{eq:risk-premium-def} with $\tilde\varepsilon=t\tilde z$. Then $\pi$ is
	twice differentiable at $t=0$ with $\pi(0)=\pi'(0)=0$ and
	\[
		\pi(t)=\frac{1}{2}\,A(w)\,t^2+o(t^2)
		\qquad\text{as }t\downarrow0,
		\qquad\text{where}\quad
		A(w)\eqdef-\frac{u''(w)}{u'(w)}
	\]
	is the \dfnas{coefficient of absolute risk aversion}{absolute risk aversion}.
	Writing $\sigma^2=t^2$ for the variance of the gamble, this is
	$\pi\approx\tfrac12A(w)\sigma^2$.
\end{theorem}

\begin{proof}
	Define $\Phi(\pi,t)\eqdef u(w-\pi)-\E[u(w+t\tilde z)]$. Boundedness of
	$\tilde z$ and continuity of $u$ and its derivatives permit differentiation under
	the expectation on a neighbourhood of $(0,0)$, by
	\autoref{thm:differentiation-under-integral}. We have $\Phi(0,0)=0$ and
	$\partial\Phi/\partial\pi=-u'(w)\neq0$ at $(0,0)$, so by the implicit function
	theorem (\autoref{thm:implicit-function}) there is a twice continuously
	differentiable $\pi(t)$ near $t=0$ with $\pi(0)=0$ and $\Phi(\pi(t),t)=0$.

	Expand the right-hand side of \eqref{eq:risk-premium-def} with the Lagrange
	remainder of \autoref{thm:taylor-1d} at order $2$: for each realisation $z$,
	\[
		u(w+tz)=u(w)+u'(w)tz+\tfrac12u''(w)t^2z^2+\tfrac16u'''(\xi_z)t^3z^3 .
	\]
	Taking expectations and using $\E[\tilde z]=0$, $\E[\tilde z^2]=1$ and
	boundedness of $\tilde z$ and of $u'''$ near $w$,
	\[
		\E\bigl[u(w+t\tilde z)\bigr]=u(w)+\tfrac12u''(w)t^2+O(t^3).
	\]
	Expanding the left-hand side at order $2$ in $\pi$,
	$u(w-\pi)=u(w)-u'(w)\pi+O(\pi^2)$. Equating and using $\pi=O(t^2)$, which
	makes $O(\pi^2)=O(t^4)$,
	\[
		-u'(w)\pi(t)+O(t^4)=\tfrac12u''(w)t^2+O(t^3),
	\]
	so $\pi(t)=-\tfrac12\bigl(u''(w)/u'(w)\bigr)t^2+O(t^3)$, which is the claim.
\end{proof}

\begin{remark}[Why the two sides are expanded to different orders]
	\label{rem:risk-premium-orders}
	The theorem explains the asymmetry that the slide derivation leaves implicit.
	Because $\pi$ is itself of order $t^2$, the term $-u'(w)\pi$ on the left is
	already of order $t^2$, and the next term $\tfrac12u''(w)\pi^2$ is of order
	$t^4$, hence negligible at the order retained. On the right the linear term
	vanishes because $\E[\tilde\varepsilon]=0$, so the leading non-trivial term is
	the quadratic one. Expanding the left side to first order and the right to
	second is therefore expanding \emph{both} to order $t^2$: the orders in $\pi$
	and in $t$ are different bookkeeping for the same expansion.
\end{remark}

\begin{remark}[The variance identity]\label{rem:variance-identity}
	The step $\E[\tilde\varepsilon^2]=\sigma^2$ uses
	$\Var[X]=\E[X^2]-\bigl(\E[X]\bigr)^2$, which follows by expanding
	$\E\bigl[(X-\E X)^2\bigr]=\E[X^2]-2\E[X]\E[X]+(\E X)^2$ and is
	\autoref{prop:variance-identity}. It is exactly here that the mean-zero
	normalisation of the gamble is used; for a gamble with non-zero mean the
	first-order term does not vanish and the premium contains a term of order $t$
	that reflects the change in expected wealth rather than risk
	\autocite[slide 42]{fukuda2026calculus}.
\end{remark}

\paragraph{Economic reading.} Three consequences. The premium is proportional to
the variance, so risk aversion of this local kind prices only second moments; a
gamble's skewness enters at order $t^3$ and is invisible in the approximation.
The factor of proportionality is $A(w)/2$, a property of preferences alone,
which is why $A$ is comparable across decision makers and across wealth levels
while $u''$ alone is not --- $u''$ changes under the affine transformations of
$u$ that leave preferences unchanged, and the ratio $-u''/u'$ does not. Finally,
the premium is second order in the size of the gamble: for small gambles a risk
averse agent is approximately risk neutral, which is the formal statement behind
the practice of ignoring risk in first-order comparative statics and the reason
that observed rejections of small-stakes gambles are hard to reconcile with
expected utility.

\section{Chapter Summary}\label{sec:c1-summary}

The derivative is the coefficient of the best affine approximation
\eqref{eq:linear-approx}, and differentiability implies continuity but not
conversely. At an interior maximiser or minimiser of a differentiable function
the derivative vanishes (\autoref{thm:fermat}); the conclusion fails at
boundaries, at kinks, and as a sufficient condition. Under concavity the
first-order condition becomes sufficient and identifies the global maximiser
(\autoref{thm:foc-concave}), because a differentiable concave function lies
below each of its tangents. Rolle's theorem and the mean value theorem convert
information about the derivative into information about the function, and are the
engine behind both the monotonicity criterion and Taylor's theorem. The
differentiation rules --- linearity, product, quotient, chain, inverse ---
are proved here in the forms actually used later, with the inverse rule carrying
the economic content that demand and inverse demand have reciprocal slopes.

Taylor's theorem states that the polynomial with coefficients
\eqref{eq:taylor-coefficients} approximates $f$ with an error given exactly by
\eqref{eq:lagrange-remainder} or asymptotically by
\autoref{cor:peano}. Applied to $\log(1+x)$ it justifies log differences as
growth rates, the Fisher equation, and log-linearisation, in each case with a
quantified second-order error. Applied to the certainty-equivalence condition it
produces the Arrow--Pratt formula $\pi\approx\tfrac12A(w)\sigma^2$, with the
orders of expansion on the two sides reconciled by the observation that the
premium is itself second order in the scale of the gamble. For the firm, the
first-order condition gives price equal to marginal cost under price taking and
the Lerner markup condition under price setting, the latter implying that a
monopolist operates only where demand is elastic.

\begin{exercise}\label{ex:c1-1}
	Let $c(q)=\tfrac12\gamma q^2$ with $\gamma>0$ and $P(q)=a-bq$ with
	$a,b>0$. Compute $q^\ast$ and $P(q^\ast)$ for both firm problems, verify
	\eqref{eq:lerner} directly, and show that the monopoly quantity is exactly half
	the competitive quantity when $\gamma=0$. Identify where
	\autoref{ass:firm} fails when $\gamma=0$ and why the conclusion survives anyway.
\end{exercise}

\begin{exercise}\label{ex:c1-2}
	Show that a function may be differentiable everywhere with discontinuous
	derivative, using $f(x)=x^2\sin(1/x)$ for $x\neq0$ and $f(0)=0$. Compute
	$f'(0)$ from the definition, compute $f'(x)$ for $x\neq0$, and show that
	$\lim_{x\to0}f'(x)$ does not exist. Explain why this does not contradict
	\autoref{thm:mvt}.
\end{exercise}

\begin{exercise}\label{ex:c1-3}
	For CRRA utility $u(c)=c^{1-\gamma}/(1-\gamma)$ with $\gamma>0$, $\gamma\neq1$,
	compute $A(w)$ and the coefficient of relative risk aversion $wA(w)$. Show that
	the risk premium for a multiplicative gamble $w(1+t\tilde z)$ is approximately
	$\tfrac12\gamma wt^2$, and contrast with \autoref{thm:arrow-pratt}. Verify that
	$\lim_{\gamma\to1}u(c)$, after subtracting the constant $1/(1-\gamma)$, equals
	$\log c$.
\end{exercise}
